% GENERATED FILE. Edit canonical lessons, then run npm run generate:books.
% canonical-source-hash: fnv1a64:443a6ef72393d715
% canonical-lessons: ZH-C481-yaojin, ZH-W481-jin-tight, ZH-C481-bingqing, ZH-C481-bingyin, ZH-C481-ganran, ZH-C481-bingdu

\chapter{Is It Serious?}
\label{ch:zh-is-it-serious}
\hlchaptermodality{481}

\begin{chapteropening}
\textbf{By the end of this chapter:} \emph{I can ask whether it is serious, talk about a patient's condition, the cause of an illness, infection and viruses.}

Complete the last lesson of chapter 481: I can ask whether it is serious, talk about a patient's condition, the cause of an illness, infection and viruses.
\end{chapteropening}

\section[\texorpdfstring{\zh{要紧} (yàojǐn)}{要紧 (yàojǐn)}]{\zh{要紧} (yàojǐn) — urgent; a cause for concern}

\label{lesson:ZH-C481-yaojin}



\begin{quote}
Before the new one: say the Mandarin for to be allergic, then the Mandarin for to lie awake at night.
\end{quote}

\subsection*{You'll want to know: \zh{要紧}}
\textbf{\zh{要紧}} — \emph{yàojǐn} — \textquotedblleft{}urgent; a cause for concern\textquotedblright{}.

\begin{quote}
\textbf{\zh{有一点儿血不要紧}} \emph{yǒu yìdiǎnr xiě bú yàojǐn} — There's a little blood, but it's nothing serious.
\end{quote}

The new character, in pieces: \zh{紧} — ten strokes in three pieces: at the top two short verticals and \zh{又} \emph{yòu}; under them two short zigzag strokes and a dot; and at the foot \zh{小} \emph{xiǎo}, three strokes. The next lesson puts it in your hand.

\begin{cousinweb}
\textbf{\zh{要紧}} (\emph{yàojǐn}) is built from \zh{要} 'to need' and \zh{紧} 'tight', both inherited from Middle Chinese, \emph{'jiewH} and \emph{kjinX} (Guangyun, in Baxter's transcription), which glosses the first 'to bind' and the second 'tightly bound'.
\end{cousinweb}

\begin{grammarlens}[title={What you've built}]
One more word for saying how it was.
\end{grammarlens}

\subsection*{Your turn}
\begin{itemize}
\raggedright
  \item \emph{Say it:} \emph{yàojǐn}
  \item \emph{Say it:} \emph{yàojǐn}, once more
  \item \emph{Say it:} \emph{yǒu yìdiǎnr xiě bú yàojǐn} — There's a little blood, but it's nothing serious.
  \item \emph{Recall:} say the Mandarin for to be allergic, then the Mandarin for to lie awake at night, then say \emph{yàojǐn} again
\end{itemize}

\begin{tcolorbox}[breakable,colback=teal!4,colframe=teal!35!black,title={Before you move on}]
What does \zh{要紧} mean? (\textquotedblleft{}Urgent; a cause for concern\textquotedblright{}.) Say it once more.
\end{tcolorbox}
\section[\texorpdfstring{\zh{紧} (jǐn)}{紧 (jǐn)}]{\zh{紧} — two verticals and \zh{又}, over a silk piece ending in \zh{小}}

\label{lesson:ZH-W481-jin-tight}



\begin{quote}
The word first, then the character inside it.

\begin{itemize}
\raggedright
  \item \emph{Recall:} say \emph{yàojǐn} — \textbf{R1}, one lesson back, before it decays
\end{itemize}

\begin{quote}
\zh{又} \zh{小} \zh{要}
\end{quote}

\emph{Read it:} each one aloud, with nothing said first — \emph{yòu}, \emph{xiǎo}, \emph{yào}
\end{quote}

\subsection*{Script — two verticals and \zh{又}, over a silk piece ending in \zh{小}}
\hlblockfigure{\detokenize{figures/ZH-W481-jin-tight-filmstrip.pdf}}{How \zh{紧} is written, stroke by stroke}

\begin{quote}
\zh{紧}
\end{quote}

\textbf{Three pieces, top to bottom.} Count the pieces, not the strokes.

\textbf{The top, four strokes.} Two short verticals side by side, then \zh{又}: a bar that turns into a stroke falling to the left, and a long stroke falling to the right.

\textbf{The middle, three strokes.} A short stroke down that turns right, a second one like it to its lower right, then a dot.

\textbf{The foot, \zh{小}, three strokes.} A vertical that hooks to the left, a dot on the left, then a dot on the right.

4 + 3 + 3 = ten.

\subsection*{Your turn}
\begin{itemize}
\raggedright
  \item \emph{Write it:} \zh{紧}, all ten, with the model in front of you
  \item \emph{Cover:} write it again from memory, and count the pieces, not the strokes
\end{itemize}

\begin{tcolorbox}[breakable,colback=teal!4,colframe=teal!35!black,title={Before you move on}]
How many pieces does \zh{紧} have? (\textbf{Three: four strokes, three and three}.) Which written characters are inside it? (\textbf{\zh{又} at the top and \zh{小} at the foot}.)

Source: \href{https://raw.githubusercontent.com/chanind/hanzi-writer-data/68d10a4b21150cae5e1ebbd223eed289cf32d90c/data/\%E7\%B4\%A7.json}{Hanzi Writer Data \zh{紧}, pinned stroke paths and medians}.
\end{tcolorbox}
\section[\texorpdfstring{\zh{病情} (bìngqíng)}{病情 (bìngqíng)}]{\zh{病情} (bìngqíng) — a patient's condition; how an illness is going}

\label{lesson:ZH-C481-bingqing}



\begin{quote}
Before the new one: say the Mandarin for to lie awake at night, then the Mandarin for serious or urgent.
\end{quote}

\subsection*{You'll want to know: \zh{病情}}
\textbf{\zh{病情}} — \emph{bìngqíng} — \textquotedblleft{}a patient's condition; how an illness is going\textquotedblright{}.

\begin{quote}
\textbf{\zh{他的病情不要紧你不要伤心}} \emph{tā de bìngqíng bú yàojǐn nǐ bú yào shāngxīn} — His condition isn't serious, so don't be upset.
\end{quote}

\begin{cousinweb}
\textbf{\zh{病情}} (\emph{bìngqíng}) is built from \zh{病} 'illness' and \zh{情} 'the state of things', both inherited from Middle Chinese, \emph{bjaengH} and \emph{dzjeng} (Guangyun, in Baxter's transcription), which glosses the first as an illness growing worse.
\end{cousinweb}

\begin{grammarlens}[title={What you've built}]
One more word for saying how it was.
\end{grammarlens}

\subsection*{Your turn}
\begin{itemize}
\raggedright
  \item \emph{Say it:} \emph{bìngqíng}
  \item \emph{Say it:} \emph{bìngqíng}, once more
  \item \emph{Say it:} \emph{tā de bìngqíng bú yàojǐn nǐ bú yào shāngxīn} — His condition isn't serious, so don't be upset.
  \item \emph{Recall:} say the Mandarin for to lie awake at night, then the Mandarin for serious or urgent, then say \emph{bìngqíng} again
\end{itemize}

\begin{tcolorbox}[breakable,colback=teal!4,colframe=teal!35!black,title={Before you move on}]
What does \zh{病情} mean? (\textquotedblleft{}A patient's condition; how an illness is going\textquotedblright{}.) Say it once more.
\end{tcolorbox}
\section[\texorpdfstring{\zh{病因} (bìngyīn)}{病因 (bìngyīn)}]{\zh{病因} (bìngyīn) — cause of an illness}

\label{lesson:ZH-C481-bingyin}



\begin{quote}
Before the new one: say the Mandarin for serious or urgent, then the Mandarin for how an illness is going.
\end{quote}

\subsection*{You'll want to know: \zh{病因}}
\textbf{\zh{病因}} — \emph{bìngyīn} — \textquotedblleft{}cause of an illness\textquotedblright{}.

\begin{quote}
\textbf{\zh{牙医看了半天还没找到病因}} \emph{yáyī kàn le bàntiān hái méi zhǎo dào bìngyīn} — The dentist has been looking for ages and still hasn't found the cause.
\end{quote}

\begin{cousinweb}
\textbf{\zh{病因}} (\emph{bìngyīn}) is built from \zh{病} 'illness' and \zh{因} 'a cause', both inherited from Middle Chinese, \emph{bjaengH} and \emph{'jin} (Guangyun, in Baxter's transcription), which glosses the second 'to rely on' and 'a cause'.
\end{cousinweb}

\begin{grammarlens}[title={What you've built}]
One more word for saying how it was.
\end{grammarlens}

\subsection*{Your turn}
\begin{itemize}
\raggedright
  \item \emph{Say it:} \emph{bìngyīn}
  \item \emph{Say it:} \emph{bìngyīn}, once more
  \item \emph{Say it:} \emph{yáyī kàn le bàntiān hái méi zhǎo dào bìngyīn} — The dentist has been looking for ages and still hasn't found the cause.
  \item \emph{Recall:} say the Mandarin for serious or urgent, then the Mandarin for how an illness is going, then say \emph{bìngyīn} again
\end{itemize}

\begin{tcolorbox}[breakable,colback=teal!4,colframe=teal!35!black,title={Before you move on}]
What does \zh{病因} mean? (\textquotedblleft{}Cause of an illness\textquotedblright{}.) Say it once more.
\end{tcolorbox}
\section[gǎnrǎn]{gǎnrǎn — to become infected}

\label{lesson:ZH-C481-ganran}



\begin{quote}
Before the new one: say the Mandarin for how an illness is going, then the Mandarin for the cause of an illness.
\end{quote}

\subsection*{You'll want to know: gǎnrǎn}
\emph{gǎnrǎn} — \textquotedblleft{}to become infected\textquotedblright{}.

\begin{quote}
\emph{wǒ gǎnjué shāngkǒu gǎnrǎn le} — I think the wound has become infected.
\end{quote}

Its characters come later; for now it is a word for the ear.

Two third tones meet in this word: gǎn is said with a rising tone before rǎn, and the spelling keeps both marks.

\begin{cousinweb}
\emph{gǎnrǎn} is built from \zh{感} 'to move' and \emph{rǎn} 'to dye', both inherited from Middle Chinese, \emph{komX} and \emph{nyemX} (Guangyun, in Baxter's transcription), which glosses the first 'to move' and the second 'to dye'. The medical sense is modern.
\end{cousinweb}

\begin{grammarlens}[title={What you've built}]
One more word for saying how it was.
\end{grammarlens}

\subsection*{Your turn}
\begin{itemize}
\raggedright
  \item \emph{Say it:} \emph{gǎnrǎn}
  \item \emph{Say it:} \emph{gǎnrǎn}, once more
  \item \emph{Say it:} \emph{wǒ gǎnjué shāngkǒu gǎnrǎn le} — I think the wound has become infected.
  \item \emph{Recall:} say the Mandarin for how an illness is going, then the Mandarin for the cause of an illness, then say \emph{gǎnrǎn} again
\end{itemize}

\begin{tcolorbox}[breakable,colback=teal!4,colframe=teal!35!black,title={Before you move on}]
What does gǎnrǎn mean? (\textquotedblleft{}To become infected\textquotedblright{}.) Say it once more.
\end{tcolorbox}
\section[bìngdú]{bìngdú — virus}

\label{lesson:ZH-C481-bingdu}



\begin{quote}
Before the new one: say the Mandarin for the cause of an illness, then the Mandarin for to become infected.
\end{quote}

\subsection*{You'll want to know: bìngdú}
\emph{bìngdú} — \textquotedblleft{}virus\textquotedblright{}.

\begin{quote}
\emph{zhège bìngdú bǐ gǎnmào kěpà} — This virus is worse than a cold.
\end{quote}

Its characters come later; for now it is a word for the ear.

\begin{cousinweb}
\emph{bìngdú} is built from \zh{病} 'illness' and \emph{dú} 'poison', both inherited from Middle Chinese, \emph{bjaengH} and \emph{dowk} (Guangyun, in Baxter's transcription), which glosses the second 'harm' and, after an older book, as a plant that harms people. The sense 'virus' is modern.
\end{cousinweb}

\begin{grammarlens}[title={What you've built}]
One more word for saying how it was.
\end{grammarlens}

\subsection*{Your turn}
\begin{itemize}
\raggedright
  \item \emph{Say it:} \emph{bìngdú}
  \item \emph{Say it:} \emph{bìngdú}, once more
  \item \emph{Say it:} \emph{zhège bìngdú bǐ gǎnmào kěpà} — This virus is worse than a cold.
  \item \emph{Recall:} say the Mandarin for the cause of an illness, then the Mandarin for to become infected, then say \emph{bìngdú} again
\end{itemize}

\begin{tcolorbox}[breakable,colback=teal!4,colframe=teal!35!black,title={Before you move on}]
What does bìngdú mean? (\textquotedblleft{}Virus\textquotedblright{}.) Say it once more.
\end{tcolorbox}
